\documentclass[11pt]{article}

% ── PACKAGES ─────────────────────────────────────────────────────────────────
\usepackage[margin=1.25in]{geometry}
\usepackage{amsmath, amssymb, amsthm}
\usepackage{enumitem}
\usepackage{booktabs}
\usepackage{hyperref}
\usepackage{natbib}
\usepackage{microtype}

% ── THEOREM ENVIRONMENTS ──────────────────────────────────────────────────────
\newtheorem{theorem}{Theorem}[section]
\newtheorem{proposition}[theorem]{Proposition}
\newtheorem{corollary}[theorem]{Corollary}
\newtheorem{definition}[theorem]{Definition}
\newtheorem{remark}[theorem]{Remark}

% ── TITLE ────────────────────────────────────────────────────────────────────

\title{When Does a Stable Winner Exist? Coherence in Linear Gaussian
Dynamical Systems with an Evolution Export Test\thanks{Preprint.
Not peer reviewed.}}

\author{Kunal Bhatia\\
Independent Researcher, Heidelberg\\
\texttt{kunal29bhatia@gmail.com}}

\date{March 2026}

% ─────────────────────────────────────────────────────────────────────────────

\begin{document}

\maketitle

% ── ABSTRACT ─────────────────────────────────────────────────────────────────

\begin{abstract}
When comparing descriptions of a dynamical system across multiple prediction
tasks, it is natural to ask whether one description is globally best ---
a stable winner over the full target family.
We give an exact answer for finite-dimensional linear Gaussian dynamical
systems (LGDS): a stable winner exists within the rank-$r$ observer class
if and only if the target family is $r$-coherent, meaning that all target
information matrices share a common top-$r$ eigenspace.
When coherence fails, no rank-$r$ observer is simultaneously optimal for
all targets and the global-privilege gap is strictly positive.
We then extract an empirical diagnostic from the theorem:
supported pairwise sign flips across a target family diagnose empirical
incoherence and motivate a prediction of winner-absence.
We test this diagnostic in an individual-based evolutionary simulation
battery spanning six regimes, four descriptions, and two task families
(predictive and causal).
In a pre-aggregated adapter analysis, all 11 classified incoherent
regime-family cells are winner-absent (0 mismatches; pre-registered
pass condition met).
Current support is strongest for the negative half of the diagnostic:
empirically incoherent cells are winner-absent.
No coherent cells were observed, so the positive half remains untested.
\end{abstract}

% ── SECTION 1: INTRODUCTION ──────────────────────────────────────────────────

\section{Introduction}

Science frequently compares descriptions of the same system.
A coarse spatial summary and a fine structural count are both legitimate
measurements of a cellular automaton;
a gene-frequency composition and a lineage-level profile are both legitimate
descriptions of an evolving population.
When predictions are made from each description, they may agree or disagree,
and one description may be more accurate than another for a particular
future quantity.

The natural next question is whether one description is \emph{globally} best:
not merely better for one target, but the best description to use regardless
of what is being predicted.
In practice this question is often answered implicitly, by convention or
by computational convenience.
The formal question --- when is a stable winner over a target family even
possible? --- has not, to our knowledge, been given a precise answer in
this form.

This paper gives one.
For finite-dimensional linear Gaussian dynamical systems, among rank-$r$
linear observers and squared-error prediction targets, a stable winner
exists if and only if the target family satisfies an algebraic condition
we call $r$-coherence: all target information matrices share a common
top-$r$ eigenspace.
When this condition fails, the loss from any fixed observer applied to
a non-native target is strictly positive, and no observer can be
simultaneously optimal for every target in the family.

The theorem is narrow by design.
It does not claim that description-dependent privilege is universal, that
all real systems are incoherent, or that any particular prior empirical
result is an instance of the theorem.
It claims only that, in the exact LGDS setting, the coherence condition
is both necessary and sufficient for the existence of a stable winner.

The theorem's operative structure --- pairwise sign stability across a
target family --- suggests an empirical diagnostic that may travel beyond
LGDS.
We formalise this diagnostic (Section~\ref{sec:diagnostic}) and test it in
an evolutionary simulation battery (Section~\ref{sec:export}).
The test is not a theorem extension; it is an empirical probe of
transportability.

Four prior empirical results motivated the formal question.
In post-transient Game of Life, the scale that best predicts a future
observable depends on which observable is chosen~\citep{bhatia2026ca}.
In kinetic Ising dynamics, fine-description advantage reverses between
ordered and disordered phases~\citep{bhatia2026ising}.
In an open Bose--Hubbard chain, a causal handle effective in one coupling
regime reverses in another~\citep{bhatia2026bh}.
In three-dimensional N-body simulations, the winning description depends
on the initial-condition family~\citep{bhatia2026nbody}.
Each result posed a winner-stability question in a different system;
none supplied a formal criterion for when a stable winner should exist.
Section~\ref{sec:prior} repositions these results as empirical pressure
for the question the theorem now answers, not as theorem-governed instances.

The paper is organised as follows.
Section~\ref{sec:framework} defines descriptions, target families, scores,
coherence, and stable winnerhood.
Section~\ref{sec:theorem} states and proves the main theorem.
Section~\ref{sec:diagnostic} extracts the exported diagnostic.
Section~\ref{sec:prior} briefly repositions the prior empirical results.
Section~\ref{sec:export} reports the evolution export test.
Section~\ref{sec:scope} states scope and limitations.
Section~\ref{sec:discussion} discusses implications.

% ── SECTION 2: FRAMEWORK ─────────────────────────────────────────────────────

\section{Framework}
\label{sec:framework}

\subsection{Descriptions and target families}

Let $\mathcal{X}$ denote the state space of a dynamical system.
A \emph{description} is a measurable function $D : \mathcal{X} \to \mathbb{R}^d$
that maps a system state to a fixed-dimensional summary.
We work with a finite family of descriptions
$\mathcal{D} = \{D_1, \dots, D_k\}$,
where each $D_i$ represents a distinct representational choice ---
a coarse-graining, an observable class, or a level of description.

A \emph{target family} is a finite set
$\mathcal{T} = \{Y_1, \dots, Y_m\}$
of scalar prediction targets,
where each $Y_j$ is a measurable function of the system state at a fixed
future horizon $\tau \geq 1$.
The target family formalises the question ``what is the description for?''

\subsection{Scores}

For each description--target pair $(D_i, Y_j)$,
a \emph{score} $s(D_i; Y_j) \in \mathbb{R}$ measures predictive performance,
with higher values indicating better performance.
The score is computed on held-out data or, where exact, from an analytic
risk formula.
The pairwise gap
\begin{equation}
    \Delta_{ij}(Y) \;=\; s(D_i;\, Y) \,-\, s(D_j;\, Y)
\end{equation}
is positive when $D_i$ outperforms $D_j$ on target $Y$.

A target $Y \in \mathcal{T}$ is \emph{estimable} for the pair $(D_i, D_j)$
if the score difference $\Delta_{ij}(Y)$ can be evaluated
(in the empirical setting, if the sample size exceeds a pre-specified floor).
A gap is \emph{supported} if the estimable gap is nonzero ---
either exactly, in the analytic case, or with a confidence interval
that excludes zero, in the empirical case.

\subsection{Coherence}

\begin{definition}[Coherence]
A target family $\mathcal{T}$ is \emph{coherent} with respect to
description family $\mathcal{D}$ if,
for every ordered pair $(D_i, D_j)$ with $i \neq j$,
the sign of $\Delta_{ij}(Y)$ is the same across all estimable
$Y \in \mathcal{T}$ where the gap is supported.

Equivalently, no pair $(D_i, D_j)$ exhibits a \emph{sign flip}:
two estimable targets $Y_a, Y_b \in \mathcal{T}$ such that
$\Delta_{ij}(Y_a) > 0$ and $\Delta_{ij}(Y_b) < 0$
with both gaps supported.
\end{definition}

Coherence is a property of the joint structure of the description family and
the target family.
It does not require a unique winner on every target;
ties may remain even when pairwise ordering is stable.

\subsection{Stable winnerhood}

\begin{definition}[Stable winner]
A description $D^* \in \mathcal{D}$ is a \emph{stable winner} over
$\mathcal{T}$ within $\mathcal{D}$ if
\begin{equation}
    s(D^*;\, Y) \;\geq\; s(D_j;\, Y)
    \quad \text{for all } D_j \neq D^* \text{ and all } Y \in \mathcal{T},
\end{equation}
with strict inequality on at least one target for each competitor.
\end{definition}

A stable winner is a description that is not dominated on any prediction target
within the family.
In the exact analytic setting, this requires simultaneous optimality.
In the empirical export section (Section~\ref{sec:export}),
stable winnerhood is operationalised by a pre-registered rule
defined separately from the theorem; that operationalisation is
intentionally distinct from the exact definition above.

% ── SECTION 3: THEOREM ───────────────────────────────────────────────────────

\section{Task-Relative Descriptive Privilege in Linear Gaussian Dynamical Systems}
\label{sec:theorem}

\subsection{System and observer family}

Let the state evolve as
\begin{equation}
    x_{t+1} = A x_t + \xi_t,
    \qquad \xi_t \sim \mathcal{N}(0, Q),
    \qquad x_t \in \mathbb{R}^n,
\end{equation}
with $A$ stable (spectral radius less than one) and $Q \succ 0$.
The unique stationary covariance $\Sigma$ satisfies $A\Sigma A^\top + Q = \Sigma$
and is strictly positive definite.

The description family is the rank-$r$ linear observer class:
each description is parametrised by a row-orthonormal matrix
$U \in \mathbb{R}^{r \times n}$, $UU^\top = I_r$, acting as
\begin{equation}
    D_U(x) = U \Sigma^{-1/2} x.
\end{equation}
Pre-whitening by $\Sigma^{-1/2}$ places all observers on a common isotropic
footing.
The observer space is the Grassmannian $\mathrm{Gr}(r,n)$.

The target family consists of squared-error prediction targets of the form
$Y_{C,\tau} = C x_{t+\tau}$,
with $C \in \mathbb{R}^{p \times n}$ and integer $\tau \geq 1$.
Performance is total mean squared prediction error.

\subsection{Target information matrices and exact risk}

Given observation $z = U\Sigma^{-1/2}x_t$,
the MMSE predictor of $Cx_{t+\tau}$ is linear by Gaussian conjugacy.
The Bayes risk is
\begin{equation}
    R(U;\, C, \tau)
    \;=\;
    \mathrm{tr}(C\Sigma C^\top)
    \;-\;
    \mathrm{tr}\!\left(U M_{C,\tau} U^\top\right),
    \label{eq:bayes-risk}
\end{equation}
where the \emph{target information matrix} is
\begin{equation}
    M_{C,\tau}
    \;=\;
    \Sigma^{1/2}(A^\tau)^\top C^\top C A^\tau \Sigma^{1/2}
    \;\in\; \mathbb{R}^{n \times n}.
    \label{eq:info-matrix}
\end{equation}
$M_{C,\tau}$ is symmetric positive semidefinite.
The formula is exact; no approximation is involved.

Setting the score $s(D_U; Y_{C,\tau}) = -R(U; C, \tau)$
(higher score = lower risk),
minimising risk over $\mathrm{Gr}(r,n)$ is equivalent to maximising
$\mathrm{tr}(U M_{C,\tau} U^\top)$.

\subsection{Top-eigenspace optimality}

\begin{proposition}[Ky Fan optimality]
\label{prop:ky-fan}
For fixed rank $r$, the observer minimising $R(U; C, \tau)$ over
$\mathrm{Gr}(r,n)$ has row space equal to the top-$r$ eigenspace of
$M_{C,\tau}$.
The minimum risk is
\begin{equation}
    R^* = \mathrm{tr}(C\Sigma C^\top) - \sum_{i=1}^{r} \lambda_i(M_{C,\tau}),
\end{equation}
where $\lambda_1 \geq \cdots \geq \lambda_n$ are the eigenvalues of
$M_{C,\tau}$.
Uniqueness (up to rotation within the eigenspace) holds when
$\lambda_r > \lambda_{r+1}$.
\end{proposition}

\begin{proof}
Minimising $R$ is equivalent to maximising
$\mathrm{tr}(U M_{C,\tau} U^\top)$ subject to $UU^\top = I_r$.
By the Ky Fan variational principle, the maximum equals
$\sum_{i=1}^r \lambda_i(M_{C,\tau})$,
achieved exactly when the row space of $U$ spans the corresponding eigenspace.
\end{proof}

\subsection{Main theorem}

The framework vocabulary of Section~\ref{sec:framework} maps onto the LGDS
setting as follows.
Each rank-$r$ observer $U$ defines a description $D_U$.
The score is $s(D_U; Y_{C,\tau}) = -R(U; C, \tau)$.
A description $D_U$ is a stable winner over $\mathcal{T}$ within the
rank-$r$ class if its row space simultaneously minimises risk for every
target in the family.

\begin{definition}[$r$-coherence]
\label{def:r-coherence}
A target family $\mathcal{T} = \{(C_k, \tau_k)\}$ is \emph{$r$-coherent}
with respect to $(A, \Sigma)$ if all information matrices
$\{M_{C_k, \tau_k}\}$ share a common top-$r$ eigenspace $V^*$.
\end{definition}

In LGDS, $r$-coherence is the structural condition that characterises
when a stable winner exists within the rank-$r$ observer class.
The proof below is a direct consequence of the Ky Fan variational
principle; the contribution is the packaging of coherence and stable
winnerhood as named conditions with an empirical diagnostic that can
be deployed outside the linear Gaussian class.

\begin{theorem}[Stable winner iff $r$-coherent]
\label{thm:main}
In a finite-dimensional LGDS at stationarity,
among rank-$r$ linear observers and squared-error prediction targets
$\mathcal{T} = \{Y_{C_k,\tau_k}\}$:
\begin{enumerate}[label=(\roman*)]
    \item \textbf{Existence.}
    A stable winner exists over $\mathcal{T}$ within the rank-$r$ observer
    class if and only if $\mathcal{T}$ is $r$-coherent.

    \item \textbf{No-global-privilege corollary.}
    If two targets in $\mathcal{T}$ have information matrices with distinct
    top-$r$ eigenspaces, no rank-$r$ observer is simultaneously optimal for
    both; no stable winner exists.

    \item \textbf{Coherent case.}
    When $\mathcal{T}$ is $r$-coherent with common eigenspace $V^*$,
    the unique stable winner (up to rotation) is the observer whose row
    space equals $V^*$.
\end{enumerate}
\end{theorem}

\begin{proof}
\emph{(i), necessity.}
Suppose $D_{U^*}$ is a stable winner.
By Proposition~\ref{prop:ky-fan}, optimality for target $Y_{C_k,\tau_k}$
requires the row space of $U^*$ to equal the top-$r$ eigenspace of
$M_{C_k,\tau_k}$.
Since $D_{U^*}$ must be optimal for every target in $\mathcal{T}$,
all information matrices share the common eigenspace $\mathrm{rowspace}(U^*)$.
Hence $\mathcal{T}$ is $r$-coherent.

\emph{(i), sufficiency.}
Suppose $\mathcal{T}$ is $r$-coherent with common eigenspace $V^*$.
Let $U^*$ be any row-orthonormal basis for $V^*$.
By Proposition~\ref{prop:ky-fan}, $D_{U^*}$ minimises risk simultaneously
for every target in $\mathcal{T}$ and is therefore a stable winner.

\emph{(ii).}
If $M_{C_1,\tau_1}$ and $M_{C_2,\tau_2}$ have distinct top-$r$ eigenspaces,
no single row space can equal both.
By Proposition~\ref{prop:ky-fan}, no rank-$r$ observer is optimal for both
targets simultaneously.

\emph{(iii).}
Follows directly from the sufficiency direction of (i) and the uniqueness
clause of Proposition~\ref{prop:ky-fan} when $\lambda_r > \lambda_{r+1}$
for each target's information matrix.
\end{proof}

\subsection{$r$-coherence is structurally demanding}

$r$-coherence requires every target in the family to be a linear functional
of the same $r$-dimensional sufficient statistic, across all prediction
horizons.
This demands strong alignment between the dynamics $(A, \Sigma)$ and the
task family $\mathcal{T}$.
It is an algebraic equality condition, not an approximate or asymptotic one.

Outside $r$-coherence, there is no rank-$r$ observer that is simultaneously
optimal for all targets.
The cost of forcing any fixed observer onto a non-native target is given
exactly by equation~\eqref{eq:bayes-risk}:
the global-privilege gap
\begin{equation}
    \Gamma
    \;=\;
    \min_{U \in \mathrm{Gr}(r,n)}
    \max_k \left[ R(U;\, C_k,\tau_k) - R^*_k \right]
    \;>\; 0
\end{equation}
measures the irreducible worst-case regret under task conflict.

\subsection{Numerical illustration}

We illustrate Theorem~\ref{thm:main} with a concrete instance
($n = 8$, $r = 2$, seed 100).\footnote{This numerical illustration
recapitulates the construction from \citet{bhatia2026lgds} for
self-containedness.}
Two targets $T_1 = ([e_1^\top; e_2^\top],\, 1)$ and
$T_2 = ([e_5^\top; e_6^\top],\, 5)$ have information matrices whose top-2
eigenspaces subtend principal angles of $54.1^\circ$ and $24.0^\circ$:
no single rank-2 observer is simultaneously optimal.
The minimax observer $U_{\mathrm{joint}}$, found by projected optimisation
over $\mathrm{Gr}(2,8)$, achieves $\Gamma = 0.026 > 0$
(Table~\ref{tab:regret}).

A structurally coherent instance is constructed by design
(normal $A_{\mathrm{sym}}$, targets aligned to the common top-2 eigenspace):
the maximum pairwise principal angle across all four targets is
$0.000^\circ$ to machine precision,
and perturbations of magnitude $\varepsilon = 0.1$ to $C$
(20 random draws) produce a maximum principal angle of $0.102$~rad ---
well below the fragility threshold of $0.15$~rad.
Structural coherence is perturbation-robust,
distinguishing it from accidental near-coincidences that break under
small perturbations.

\begin{table}[t]
\centering
\caption{Regret matrix, no-global-privilege instance ($n=8$, $r=2$,
seed 100). All values from the exact analytic Bayes risk formula.}
\label{tab:regret}
\begin{tabular}{lcc}
\toprule
Observer & Regret on $T_1$ & Regret on $T_2$ \\
\midrule
$U^*_1$ (optimal for $T_1$) & 0.000 & 0.036 \\
$U^*_2$ (optimal for $T_2$) & 0.763 & 0.000 \\
$U_{\mathrm{joint}}$ (minimax) & 0.026 & 0.026 \\
\bottomrule
\end{tabular}
\end{table}

% ── SECTION 4: EXPORTED DIAGNOSTIC ──────────────────────────────────────────

\section{The Exported Diagnostic}
\label{sec:diagnostic}

\subsection{From eigenspaces to sign flips}

Theorem~\ref{thm:main} is exact within LGDS.
This section extracts from it an \emph{empirical diagnostic} that can be
evaluated in systems that are not LGDS.

The theorem's operative condition --- whether information matrices share a
common top-$r$ eigenspace --- is not directly testable in systems where the
dynamics and noise structure are unknown or nonlinear.
What \emph{is} testable in any system is the sign of pairwise score
differences across targets.

The diagnostic rests on the following observation.
In LGDS, when two targets have distinct top-$r$ eigenspaces, no stable
winner exists (Theorem~\ref{thm:main}(ii)).
By the exact risk formula~\eqref{eq:bayes-risk}, the optimal observer for
target $k$ incurs strictly positive regret on any other target with a
distinct eigenspace.
This eigenspace separation implies that no single observer is simultaneously
optimal for both targets.
In finite observer families, this loss of common optimality is naturally
probed by pairwise score-gap sign changes across targets.

Conversely, if the target family is $r$-coherent, the same observer is
optimal for every target in the rank-$r$ class;
this is the exact LGDS case in which no task conflict remains.

\subsection{The exported prediction}

This motivates a diagnostic that travels outside LGDS:

\begin{quote}
\emph{If at least one pairwise sign flip is empirically supported across a
target family, no stable family winner should exist.}
\end{quote}

This is a prediction, not a theorem.
Outside LGDS, sign flips are sufficient evidence of incoherence in the
general sense of Section~\ref{sec:framework}, but the connection to stable
winnerhood is no longer guaranteed by Theorem~\ref{thm:main}.
The export test asks whether the diagnostic nevertheless has empirical bite.

\subsection{Operationalisation}

To test the exported prediction, we require operational definitions of
sign-flip support and stable winnerhood that are independent of any
parametric model.

\paragraph{Supported sign flip.}
For an ordered pair of descriptions $(D_i, D_j)$ and an estimable target
$Y \in \mathcal{T}$, the gap $\Delta_{ij}(Y)$ is \emph{supported in
direction} $+$ if its 95\% bootstrap confidence interval excludes zero from
below ($\mathrm{lo} > 0$), and \emph{supported in direction} $-$ if the
interval excludes zero from above ($\mathrm{hi} < 0$).
A pair $(D_i, D_j)$ exhibits a \emph{supported sign flip} across $\mathcal{T}$
if there exist two estimable targets $Y_a, Y_b$ with opposite supported
directions.
A target family is \emph{empirically incoherent} if at least one pair
exhibits a supported sign flip; \emph{empirically coherent} if every pair
is sign-stable (all supported gaps point the same direction);
and \emph{unresolved} otherwise.

\paragraph{Stable family winner (empirical).}
Description $D^*$ is a \emph{stable family winner} if, for every competitor
$D_j \neq D^*$:
\begin{enumerate}[label=(\roman*)]
    \item $D^*$ beats $D_j$ on at least 75\% of estimable targets
          with CI-supported gaps;
    \item no competitor beats $D^*$ on more than 25\% of estimable targets
          with CI-supported gaps.
\end{enumerate}
Otherwise the cell is \emph{winner-absent}.

These thresholds were fixed before examination of the regime-family
results.

\paragraph{Correspondence test.}
The exported diagnostic is evaluated by the \emph{correspondence test}:
across pre-specified regime-family cells, does empirical incoherence predict
winner-absence, and does empirical coherence predict stable winnerhood?
The primary exported claim is the negative half:
incoherence $\to$ winner-absent.
The positive half (coherence $\to$ stable winner exists) requires coherent
cells to be present in the battery; their presence or absence is itself
a data outcome.

% ── SECTION 5: PRIOR ODD PAPERS ──────────────────────────────────────────────

\section{Relation to Prior Empirical Results}
\label{sec:prior}

The empirical results motivating this paper were obtained before the formal
framework above was developed.
We briefly reposition them here as instances of coherence failure, not as
evidence for a stronger unified law.

\paragraph{Cellular automata (P1, P4).}
In post-transient Conway's Game of Life, a fine component-count observer and
a coarse occupied-block observer disagree on net change over the same run.
The disagreement is predictable from early fine-scale structure
($r = -0.848$, 95\% CI $[-0.862, -0.830]$),
and the scale that best predicts a future observable depends on which
observable is chosen:
future live-cell count peaks at block size $B = 2$, while future
occupied-block count peaks at $B = 8$, with non-overlapping bootstrap
confidence intervals~\citep{bhatia2026ca}.
In the framework of Section~\ref{sec:framework}, the target family
$\{\bar{N}_\mathrm{live}, \bar{N}^{B=8}_\mathrm{occ}\}$ is empirically
incoherent: the pairwise ordering between fine and coarse descriptions
flips across the two targets.
A stable winner is not observed over this family;
the present theorem provides the formal vocabulary for understanding
that kind of failure in the exact LGDS setting.

\paragraph{Kinetic Ising model (P2).}
Fine descriptions based on locally exposed defects outperform a coarse
block-variance description for predicting future domain-wall change
($\Delta R^2 = +0.415$, 95\% CI $[+0.389, +0.443]$) in the ordered phase,
but the advantage reverses in the disordered phase
($\Delta R^2 = -0.064$)~\citep{bhatia2026ising}.
The regime dependence shows that no regime-independent stable winner exists
over the combined family of ordered and disordered targets.

\paragraph{Bose--Hubbard chain (P3).}
In an open Bose--Hubbard chain with Lindblad dephasing,
targeted additional dephasing at high-variance sites produces more local
occupation loss than matched-budget random targeting in the
$J/U = 0.30$--$0.40$ regime, but the effect reverses at $J/U = 0.12$~\citep{bhatia2026bh}.
This is a causal instance of target-relative privilege.
Because the causal axis lies outside the scope of the present theorem,
we cite it here only as empirical pressure for a general winner-instability
question, not as a theorem-governed instance.

\paragraph{N-body dynamics (P5).}
In three-dimensional self-gravitating N-body simulations, the coarse density
variance robustly predicts future clustering for bimodal initial conditions
($r = -0.975$, 95\% CI $[-0.99, -0.95]$), but the kinematic fine observable
(local velocity dispersion) shows a suggestive advantage for concentrated
profiles at small softening~\citep{bhatia2026nbody}.
The results suggest that winner identity may depend on initial-condition
family, though the fine-side concentrated-profile signal remains unresolved.
We cite this paper as pressure for the formal question, not as a settled
theorem instance.

These four cases are not re-proved here.
They are cited as recurring empirical pressure for the formal question the
theorem addresses: when is a stable winner over a target family even possible?
The theorem gives the answer for LGDS;
the export test in Section~\ref{sec:export} probes whether the answer
generalises.
Section~\ref{sec:prior} is deliberately subordinate to the theorem and
export core; its purpose is motivation, not reinterpretation.

% ── SECTION 6: EVOLUTION EXPORT TEST ─────────────────────────────────────────

\section{Evolution Export Test}
\label{sec:export}

\subsection{Setup}

We test the exported diagnostic in an individual-based evolutionary
simulation.
Four descriptions are compared:
gene-frequency composition ($D_\mathrm{gene}$),
lineage-level structure ($D_\mathrm{lin}$),
organism-level state ($D_\mathrm{org}$),
and an organism-plus-ecology composite ($D_\mathrm{eco}$).
The target family consists of four evolutionary outcomes:
persistence, transmission, dominance, and ecological state.
Two task families are evaluated separately:
predictive (5-fold cross-validated $R^2$) and causal
(mean outcome difference under matched interventions).

Six evolutionary regimes are tested:
selected Wright-Fisher (\texttt{selected\_wf}),
Moran (\texttt{moran}),
frequency-dependent selection (\texttt{freq\_dep}),
eco-evolutionary coupling (\texttt{eco\_evol}),
group-structured (\texttt{group\_structured}),
and neutral Wright-Fisher (\texttt{neutral\_wf}).
This yields 12 regime-family cells.

Full simulation parameters and the pre-registered battery design are
reported in~\citet{bhatia2026evolution}.

\subsection{Coherence and winner classification}

For each regime-family cell, pairwise gaps $\Delta_{ij}(Y)$ are computed
from pre-aggregated scores using the independence approximation:
\begin{equation}
    \widehat{\mathrm{se}}_{ij}
    \;=\;
    \sqrt{\widehat{\mathrm{se}}_i^2 + \widehat{\mathrm{se}}_j^2},
\end{equation}
where each $\widehat{\mathrm{se}}_k$ is derived from the individual 95\%
bootstrap confidence interval as
$\widehat{\mathrm{se}}_k = (\mathrm{ci\_hi}_k - \mathrm{ci\_lo}_k) /
(2 z_{0.975})$.
With $n = 100{,}000$ replicates per cell, the resulting intervals are
narrow; however, the deviation from joint replicate-level bootstrapping
remains a documented departure from the pre-specified protocol.

Coherence and winner labels are assigned by the rules in
Section~\ref{sec:diagnostic}.%
\footnote{The two classifiers consume a shared, deterministically seeded
gap table; neither classifier re-randomises the same statistic, preventing
Monte Carlo noise from producing artificial classification discrepancies.}

\subsection{Results}

Table~\ref{tab:coherence} reports coherence and winner labels for all
12 regime-family cells.
In this pre-aggregated adapter analysis,
eleven cells received a non-unresolved coherence label;
one cell (\texttt{selected\_wf}/causal) was unresolved and
excluded from the correspondence test per the pre-registered rule.

Of the 11 classified cells:
\begin{itemize}
    \item all 11 are empirically incoherent;
    \item all 11 are winner-absent;
    \item the correspondence test yields 0 mismatches.
\end{itemize}

\paragraph{Note on winner classification.}
In the current results file, organism-plus-ecology data is present only
for the \texttt{eco\_evol}/predictive cell; all other cells lack the
fourth description entirely.
Because the winner rule requires all four descriptions to be estimable,
winner-absent in 10 of 11 classified cells is a default rather than a
competitive outcome.
The exception is \texttt{eco\_evol}/predictive, where a competitive
check ran across the four descriptions and found no winner.
The incoherence result is unaffected: sign flips are established
from the three descriptions that are present across all regimes.

Given this coverage gap, the adapter run is directionally consistent
with the winner-absence prediction on the covered cells, but does not
constitute a strong transport test.
The correspondence test returns a formal pass under the locked
classification rule; that verdict should be read as
adapter-consistent rather than strongly confirmed.

The dominant sign flip is the pair $(D_\mathrm{gene}, D_\mathrm{lin})$,
which appears in 9 of 11 classified cells across both task families.
This pair flips direction across the persistence/transmission divide:
gene-frequency descriptions are more predictive of transmission outcomes
in several regimes, while lineage descriptions are more predictive of
persistence outcomes, with bootstrap-supported gaps in both directions.

No coherent cells were observed in this battery.
The positive half of the exported prediction --- that coherence should imply
stable winnerhood --- therefore remains untested.

\begin{table}[t]
\centering
\caption{Coherence and winner classification for all 12 regime-family cells.
Classified cells: 11. Unresolved: 1. Mismatches: 0.}
\label{tab:coherence}
\begin{tabular}{llllc}
\toprule
Regime & Family & Coherence & Winner & Match \\
\midrule
\texttt{selected\_wf}     & predictive & incoherent & winner-absent & \checkmark \\
\texttt{selected\_wf}     & causal     & unresolved & (excluded)    & --- \\
\texttt{moran}            & predictive & incoherent & winner-absent & \checkmark \\
\texttt{moran}            & causal     & incoherent & winner-absent & \checkmark \\
\texttt{freq\_dep}        & predictive & incoherent & winner-absent & \checkmark \\
\texttt{freq\_dep}        & causal     & incoherent & winner-absent & \checkmark \\
\texttt{eco\_evol}        & predictive & incoherent & winner-absent & \checkmark \\
\texttt{eco\_evol}        & causal     & incoherent & winner-absent & \checkmark \\
\texttt{group\_structured}& predictive & incoherent & winner-absent & \checkmark \\
\texttt{group\_structured}& causal     & incoherent & winner-absent & \checkmark \\
\texttt{neutral\_wf}      & predictive & incoherent & winner-absent & \checkmark \\
\texttt{neutral\_wf}      & causal     & incoherent & winner-absent & \checkmark \\
\bottomrule
\end{tabular}
\end{table}

\subsection{Data coverage note}

After correcting the \texttt{org\_eco} $\to$ \texttt{organism\_ecology}
name mapping, organism-plus-ecology data is present in the results file
only for the \texttt{eco\_evol}/predictive cell; all other regimes and the
full causal family have no organism-plus-ecology data.
Because the pre-registered winner rule requires all four descriptions to
be estimable on a target before that target counts toward the winner
criterion, winner classification is competitive only for
\texttt{eco\_evol}/predictive (where a competitive check ran and found no
winner) and vacuous (returning winner-absent by default) for all remaining
cells.
The incoherence classification is unaffected: the
$(D_\mathrm{gene}, D_\mathrm{lin})$ sign flip is established from the
three descriptions that are present across all regimes.
A full four-description rerun is required before winner classification
can be treated as a genuine competitive result across the battery.

\subsection{What is and is not claimed}

The evolution export test supports the following claim:

\begin{quote}
In the pre-aggregated adapter analysis of the tested evolutionary simulation
battery, the coherence diagnostic derived from Theorem~\ref{thm:main}
correctly predicts winner-absence in all classified incoherent
regime-family cells (11 of 11; 0 mismatches).
\end{quote}

It does not support:
\begin{itemize}
    \item that the LGDS theorem is proved outside LGDS;
    \item that evolutionary systems are approximately linear-Gaussian;
    \item that the positive half of the diagnostic (coherence $\to$ stable
          winner) holds in this system, since no coherent cells were observed;
    \item that winner-absence has been competitively established across all
          four descriptions in all classified cells — only
          \texttt{eco\_evol}/predictive ran a full four-description
          competitive check; all other winner-absent labels are defaults
          due to missing organism-plus-ecology data;
    \item that the result generalises beyond the six tested regimes or the
          four tested descriptions.
\end{itemize}

% ── SECTION 7: SCOPE AND LIMITATIONS ─────────────────────────────────────────

\section{Scope and Limitations}
\label{sec:scope}

\paragraph{Exactness of the theorem.}
Theorem~\ref{thm:main} is exact within finite-dimensional linear Gaussian
dynamical systems at stationarity, among rank-$r$ linear observers and
squared-error prediction targets of the form $Y = Cx_{t+\tau}$.
The proof is algebraic and involves no approximation.
The scope is narrow by design: the theorem is a precise answer to a precise
question, not a universal law about description-dependent privilege.

\paragraph{What the theorem does not cover.}
The result does not apply to nonlinear systems, non-Gaussian noise,
nonlinear observers, or prediction tasks with non-quadratic loss.
The causal axis --- whether the observer optimal for prediction differs from
the observer optimal for intervention --- is not addressed by the theorem;
it requires a separate formal treatment.
Temporal asymmetry and the direction of optimal description are also outside
the theorem's scope.

\paragraph{Status of the export test.}
The evolution export test (Section~\ref{sec:export}) asks whether the
coherence diagnostic has empirical bite outside LGDS.
Current support is strongest for the negative half:
all 11 classified incoherent cells are winner-absent (0 mismatches).
The positive half --- that coherence should imply stable winnerhood ---
remains untested because no coherent cells appeared in the current battery.
The diagnostic travels in the tested direction; the full bidirectional claim
is not yet earned.

The adapter result is subject to two documented limitations.
First, pairwise gap CIs are computed via the independence approximation
from pre-aggregated scores rather than joint replicate-level bootstrapping
as pre-specified; this is a registered departure.
Second, the current adapter result is limited by incomplete four-description
coverage in the results file.
After correcting the \texttt{org\_eco} name mapping, organism-plus-ecology
data is present only for the \texttt{eco\_evol}/predictive cell;
all other regimes and the full causal family lack that fourth description
entirely.
The residual \texttt{ecology} entry in \texttt{interv\_class} is a
separate column-leak issue.
A full four-description rerun is required before winner classification can
be treated as a genuine competitive result across the battery.
Neither limitation currently appears likely to change the qualitative
verdict, but both must be resolved before the result is treated as final.

\paragraph{Generality of the diagnostic.}
The sign-flip diagnostic is motivated by the LGDS theorem but is not
proved to characterise stable winnerhood in any other system class.
Its application to evolutionary simulations is an empirical test of
transportability, not a theorem extension.
Positive results in further system classes would strengthen the case for
the diagnostic as a general tool; null results would narrow its scope.
The program of mapping which system classes admit or fail the diagnostic
is open.

\paragraph{Relation to prior formalisms.}
The framework in Section~\ref{sec:framework} is related to, but distinct
from, several prior formalisms.
The renormalisation group and the Wilson--Kadanoff coarse-graining program
ask which features are preserved under scale transformations;
the present framework asks which description is optimal for a given
prediction task, a question that is orthogonal to preservation under
transformation.
Hoel et al.'s causal emergence asks when macro descriptions carry more
effective information than micro descriptions;
the present result is complementary but narrower, characterising the
structural conditions under which a single description is globally privileged
across a task family rather than whether macro beats micro on a fixed task.
Israeli and Goldenfeld's computational irreducibility result establishes
that some cellular automata resist exact coarse-graining~\citep{israeli2004};
the incommensurability of the two observers in the CA results is consistent
with that finding.

% ── SECTION 8: DISCUSSION ────────────────────────────────────────────────────

\section{Discussion}
\label{sec:discussion}

\paragraph{What changes in how we read the prior results.}
The four prior empirical results cited in Section~\ref{sec:prior} each
reported that winner identity depended on the target or the regime.
Without a formal criterion, these findings could be read as isolated
system-specific curiosities rather than instances of a structural
phenomenon.
The theorem sharpens that reading --- not by proving those results are LGDS
instances, but by supplying the vocabulary and the criterion that make
the pattern recognisable.
A result in which fine-description advantage reverses between phases is now
legible as a candidate instance of coherence failure over a joint
regime-indexed target family: the same description does not remain best
once the target family is enlarged across regimes.
The theorem does not prove that these systems instantiate the LGDS mechanism;
it provides the formal vocabulary for stating the winner-stability question
precisely and, in the LGDS class, answering it exactly.

\paragraph{Open problems.}
Four problems remain open after this paper.

First, the positive half of the exported diagnostic is untested.
The evolution battery produced no coherent regime-family cells,
so the prediction that coherence implies stable winnerhood was never
challenged.
A battery designed to include coherent cells --- either by construction in
a synthetic system or by identifying regimes where one description is
known to dominate --- is needed before the bidirectional claim can be
assessed.

Second, the causal axis remains untreated.
The theorem concerns prediction under the squared-error criterion;
it says nothing about which description is optimal for intervention.
The prior results include causal instances (Bose--Hubbard, evolution)
where the best predictive description and the best causal handle differ.
A formal treatment of causal-axis coherence, if it exists, would require
a different object --- likely a response-theoretic functional rather than
an information matrix.

Third, the diagnostic's scope outside LGDS is unknown.
The evolution export test supports the negative half in one system class.
Whether the same diagnostic works in continuous-field systems,
in systems with non-quadratic loss, or in systems with non-stationary
dynamics is an open empirical question.
Null results in further system classes would be as informative as positive
ones: they would identify the boundary conditions under which coherence
structure predicts winner-stability.

Fourth, a formal generalisation of the theorem is open.
The LGDS result is exact because the Bayes risk formula is exact and the
Ky Fan variational principle applies directly.
Analogous exact results in nonlinear or non-Gaussian systems would require
either a different risk characterisation or an approximate treatment.
Whether the coherence condition has a natural extension beyond the linear
Gaussian class --- and whether that extension recovers the LGDS theorem as
a special case --- is the formal problem the current result points toward.

% ── BIBLIOGRAPHY ─────────────────────────────────────────────────────────────

\bibliographystyle{plainnat}
\bibliography{lgds_refs}

\end{document}
